\section{Setting and related work}\label{sec:setting}

A \emph{block} consists of variables $x\in\R^n$, a nonempty compact domain
$D\subset\R^n$ and continuous features $F=(f_1,\dots,f_p):D\to\R^p$. In a
model, the features are expressions of the same variables, and $D$ is the
set on which they must be relaxed, typically a box intersected with the
linear rows that involve only these variables. The \emph{joint graph} of $F$
on $D$ and its convex hull are
\begin{equation}\label{eq:joint-hull}
G_D(F)=\{(x,F(x)):x\in D\},\qquad H_D(F)=\conv G_D(F),
\end{equation}
where the coordinates $w_j$ of $(x,w)\in H_D(F)$ stand for the expressions
$f_j(x)$. For a direction $(a,b)\in\R^n\times\R^p$ the (lower) \emph{support
value} is
\begin{equation}\label{eq:support}
h_D(a,b)=\min_{x\in D}\{a\T x+b\T F(x)\}.
\end{equation}

\begin{proposition}[Support description]\label{prop:support}
For every $(a,b)$ and every $\beta\le h_D(a,b)$, the inequality
$a\T x+b\T w\ge\beta$ is valid for $H_D(F)$. Moreover $H_D(F)$ is the
intersection of the half-spaces $a\T x+b\T w\ge h_D(a,b)$ over all $(a,b)$.
\end{proposition}

\begin{proof}
The inequality holds on $G_D(F)$ by definition and is preserved by convex
combinations. $G_D(F)$ is compact, so $H_D(F)$ is compact and convex, and a
point outside it can be strictly separated by a linear functional; written
as a lower bound, the separating inequality is one of the stated
half-spaces.
\end{proof}

We call $a\T x+b\T w\ge\beta$ a \emph{support cut}, and a \emph{support
inequality} when $\beta=h_D(a,b)$. Support inequalities are as strong as
possible for the block: no valid inequality in these coordinates is missing
from the family. They also combine information that separate
relaxations discard. If each $f_j$ is relaxed on its own, the result is the
intersection of the hulls of the graphs of $(x,f_j)$, which can be much larger
than $H_D(F)$ when the $f_j$ share variables (Example~\ref{ex:two-rows}). For
$n=p=1$, the family recovers the convex and concave envelopes of a single
function. If integer variables occur in a block, we drop their integrality
and keep their bounds; support cuts remain valid for the mixed-integer graph,
but they need not describe its convex hull.


\subsection{Related work}\label{sec:related}

\paragraph{Factorable and simultaneous relaxations.}
Global solvers relax nonconvex models by decomposing expressions into
elementary operations and relaxing each with an auxiliary variable
\citep{McCormick1976,SmithPantelides1999,TawarmalaniSahinidis2004,BelottiEtAl2009,BestuzhevaEtAl2025}.
The resulting relaxations ignore the dependence among expressions that share
variables, and they ignore linear constraints that link those variables
\citep{ZhuHeTawarmalani2026}. Simultaneous convexification addresses the
first point. \citet{Tawarmalani2010} develops inclusion certificates for the
convex hull of several functions of the same variables;
\citet{Ballerstein2013}, as stated by \citet[Proposition~1]{LiersEtAl2021},
characterizes the hull of a vector of functions by linear combinations, and
\citet{LiersEtAl2021} separate over it in a gas network application.
\citet{HeTawarmalani2021,HeTawarmalani2022,HeTawarmalani2024} build composite
relaxations with vectors of outer functions, \citet{ZhuHeTawarmalani2026}
treat simultaneous factorable graphs with coupling constraints, and
\citet{DavarniaRichardTawarmalani2017} convexify vectors of bilinear
functions simultaneously over polytopes; \citet{HeLiuTawarmalani2025}
describe simultaneous hulls of vectors of univariate functions through moment
hulls. For multilinear and bilinear terms,
the gap between termwise and joint relaxations is well studied
\citep{Rikun1997,LuedtkeNamazifarLinderoth2012,BolandEtAl2017,BaoKhajaviradSahinidisTawarmalani2015}.
Several bilinear and quadratic terms are relaxed jointly by the multiterm
relaxations of \citet{BaoSahinidisTawarmalani2009} and by the edge-concave
relaxations of \citet{MisenerFloudas2012}. Reduced-space relaxations avoid
auxiliary variables altogether
\citep{TsoukalasMitsos2014,BongartzMitsos2017,NajmanBongartzMitsos2021}. Our
support cuts are an instance of simultaneous convexification; what we add
concerns which blocks are worth relaxing jointly, which support problems can
be solved exactly, how the cuts enter the original model, and how they are
certified.

\paragraph{Lagrangian and aggregation cuts.}
Cuts in the original variables obtained by minimizing a weighted sum of
constraint functions over a block are Lagrangian cuts. They appear in block
decomposition methods for MINLP \citep{Nowak2005,KaruppiahGrossmann2008,WuMutsNowakHendrix2025},
in range reduction \citep{TawarmalaniSahinidis2004,GleixnerEtAl2017} and in
rigorous filtering \citep{DomesNeumaier2016}. The relaxation they describe,
convexification in the joint space of variables and constraint values, is the
primal form of the Lagrangian dual
\citep{Falk1969,Geoffrion1974,FeltenmarkKiwiel2000,LemarechalRenaud2001}, and
the set version of this statement for Lagrangian cuts is given by
\citet{ChenLuedtke2022}. Aggregations of quadratic constraints that describe
convex hulls under additional hypotheses \citep{DeyMunozSerrano2022,BlekhermanDeySun2024},
surrogate relaxations that keep an aggregated constraint nonconvex
\citep{MullerEtAl2022}, and the joint-range inequalities of
\citet{XuPokutta2026}, which convexify the joint range of two quadratic
functions and lift the result to sparse cuts, are different objects; Section~\ref{sec:feasible-hull}
relates them to our cuts.

\paragraph{Exact low-dimensional and structured quadratic hulls.}
\citet{AnstreicherBurer2010} show that the convex hull of $\{(x,xx\T)\}$ is
described by the semidefinite constraints and the constraints of the
reformulation-linearization technique (RLT) \citep{Shor1987,SheraliAdams1990}
over a simplex of dimension at most three and over a box in dimension two,
and that polytopes of dimension at most three admit an extended formulation
built from a triangulation; for the box in dimension three these constraints
are not enough \citep{AnstreicherBurer2010,BurerLetchford2009,Anstreicher2012}.
Explicit envelopes through polyhedral subdivisions are given by
\citet{TawarmalaniRichardXiong2013}.
Convex envelopes of bivariate functions over general polygons, with
supporting hyperplanes, are computed by \citet{LocatelliSchoen2014}.
Minimizing a quadratic over a polytope of fixed dimension by enumerating
faces is classical \citep[Section~2.9]{Murty1997}, quadratic programming is in
NP \citep{Vavasis1990}, and box-constrained quadratic programming on forests
is strongly polynomial \citep{DelPiaKhajavirad2026}. Explicit hulls for
star-shaped sign patterns are given by \citet{DeyKhajavirad2025} and
\citet{Khajavirad2026}, and parametric dynamic programs over trees solve
convex quadratic problems with indicators \citep{BhathenaEtAl2025}. The
gluing of local relaxations along a sparsity pattern is discussed in
Section~\ref{sec:composition-literature}.

\paragraph{Rigorous relaxations and safe cuts.}
Safe bounds and cuts for linear and mixed-integer programs correct
floating-point results with variable bounds and directed rounding
\citep{NeumaierShcherbina2004,CookEtAl2009SafeCuts,EiflerGleixner2024}; exact
MIP solvers and VIPR certificates check complete solves against the original
instance \citep{CookKochSteffyWolter2013,CheungGleixnerSteffy2017,EiflerGleixner2023},
SCIP~10 offers an exact mode for MILP \citep{SCIP10}, and
\citet{HalbigEtAl2024} compute optimality certificates for convex MINLPs. For nonlinear
problems, rigorous linear relaxations and safe linear estimators are
established \citep{BorradaileVanHentenryck2005,Kearfott2011,DomesNeumaier2012,NininMessineHansen2015},
as are verified Bernstein bounds
\citep{GarloffJanssonSmith2003JCAM,GarloffSmith2008,MunozNarkawicz2013} and
ball arithmetic \citep{Johansson2017arb}. That rounding in the problem
formulation itself must be controlled was pointed out by
\citet[Section~20]{Neumaier2004} and \citet{SchichlNeumaier2005}.
\citet{LiersEtAl2021} mention safe rounding of coefficients without details.
Cut generators are usually tested for validity against known feasible
solutions \citep{Margot2009}; replay of exact certificates is a stronger test
for the cuts it covers. SCIP relaxes its own nonlinear cuts with variable
bounds before storing them, in floating point, to protect them against the
solver's coefficient rounding \citep[Section~4.2.10]{SCIP8report}. To our
knowledge, no published MINLP implementation certifies the exported rows of
joint support cuts against the source model and the row stored by the solver.

\paragraph{Separation and oracles.}
Weak separation and weak optimization are polynomially equivalent for convex
bodies given by oracles \citep{GrotschelLovaszSchrijver1981,GrotschelLovaszSchrijver1988}.
Gilbert's nearest-point method gives an oracle procedure that returns either
a nearby point of the set or a separating hyperplane
\citep{Gilbert1966,BrierleyNavascuesVertesi2016}, and the local-cut and
Fenchel-cut schemes separate by column generation over an optimization
oracle, in exact arithmetic when needed \citep{Boyd1994,ChvatalCookEspinoza2013};
in direction space this column generation is Kelley's cutting-plane method
\citep{CheneyGoldstein1959,Kelley1960}. Oracle-based cuts for polynomial
optimization are studied by \citet{BienstockChenMunoz2020}.
Appendix~\ref{app:separation} states the corresponding finite procedure for
joint graphs and its certificate.

\paragraph{Native SCIP.}
SCIP's nonlinear constraint handler detects quadratic, bilinear, convex,
concave and other structures, uses projections of linear rows for bilinear
terms \citep{MullerSerranoGleixner2020}, and separates RLT cuts
\citep{BestuzhevaGleixnerAchterberg2025} and semidefinite minor cuts
\citep{BestuzhevaEtAl2025}. If a variable with finite bounds and
no objective coefficient occurs in a single polynomial constraint that is
concave (convex) in it, with an infinite left-hand (right-hand) side, some
optimal solution has the variable at one of its bounds; SCIP's presolve then
makes the variable binary when its bounds are $[0,1]$, or adds a bound
disjunction \citep[Section~4.2.7]{SCIP8report}, an idea that goes back to
\citet{HansenEtAl1993}. Several cut families that target the same joint
quadratic structure as our cuts are implemented but disabled by default
because they have not been found to pay off in general
\citep{BestuzhevaEtAl2025,SCIP10}: edge-concave cuts on aggregated quadratic
terms \citep{MisenerFloudas2012}, intersection cuts for quadratic constraints
from maximal quadratic-free sets
\citep{MunozSerrano2022,ChmielaMunozSerrano2023}, interminor cuts, which are
intersection cuts for vanishing $2\times2$ minors of the matrix of products
\citep[Section~4.11]{SCIP8report}, related to the outer-product-free sets of
\citet{BienstockChenMunoz2020}, and RLT cuts for hidden products
\citep{BestuzhevaGleixnerAchterberg2025}. A new separator therefore has to be
evaluated against SCIP's default and against these separators, not against
termwise McCormick relaxations; Section~\ref{sec:computations} does both.
